\documentclass[12pt,
	english,			% idioma adicional para hifenização
	french,				% idioma adicional para hifenização
	spanish,			% idioma adicional para hifenização
	brazil,				% o último idioma é o principal do documento
	]{article}

%% Language and font encodings
\usepackage[brazil]{babel}
\usepackage[utf8x]{inputenc}
\usepackage[T1]{fontenc}
\usepackage{indentfirst}
\usepackage{float}

%% Sets page size and margins
\usepackage[a4paper,top=3cm,bottom=2cm,left=3cm,right=3cm,marginparwidth=1.75cm]{geometry}

%% Useful packages
\usepackage{amsmath}
\usepackage{graphicx}
\usepackage[colorinlistoftodos]{todonotes}
\usepackage[colorlinks=true, allcolors=blue]{hyperref}
\usepackage{url}

%% Document Header
\title{Relatório I\\
\Large{Física Experimental IV}\\
\Large{Instituto de Física Gleb Wataghin, Unicamp}}

\author{Prof. Dr. Flávio Caldas da Cruz\\
Caroline Guimarães 155006\\
Lucas Rodrigues Contador 156406\\
Giovanne Lucas Dias Pereira Mariano 173317}

%% Document Begin
\begin{document}
\maketitle

\begin{abstract}

% Seu resumo do relatório, resumo não é introdução! Por exemplo, o excerto abaixo foi extraído do relatório de alunos do 1s2017, note que ele resume as atividades desenvolvidas e apresentadas no relatório, porém não faz uma introdução exagerada ao assunto abordado.
%"\textit{Calculamos a partir da montagem de circuitos RC e RLC a resistência interna do osciloscópio, bem como os valores de capacitância dos componentes utilizados. Analisamos as curvas dos circuitos passa-altas, passa-baixas e passa banda, juntamente com seus respectivos diagramas de bode, comparando-as com curvas esperadas fornecidas na literatura sobre o assunto. Utilizando um método baseado em assíndotas, estimamos valores para as frequências de ressonância e de corte, comparando-os com valores experimentais e teóricos}"

\end{abstract}

\section{Introdução}

%Descrever o propósito do experimento, as motivações de executá-los e possíveis conceitos físicos que serão explorados.

\section{Metodologia}

%Listar equipamentos (modelo e número de identificação quando possível) e componentes (valores nominais) utilizados.

\section{Resultados}

%Apresentar os resultados experimentais, inclua eventuais medidas de caracterização de componente (e.g., resistências, resistência parasítica do indutor). Avalie se é melhor apresentar o dado na forma de tabela ou gráfico. Note que um diagrama de Bode não deve ser apresentado na forma de tabela, enquanto que os valores de resistência utilizados poderiam utilizar este formato de apresentação. 
%\begin{itemize}
%\item Caso apresente dados que serão usados para ajustar um modelo teórico, aproveite o espaço do gráfico para apresentar o ajuste. Faça uma breve discussão do ajuste,incluindo os parâmetros extraídos e equações utilizadas (evite deduções nesta parte). 
%\item Sendo seu experimento constituído de mais de um conjunto de dados, utilize subseções.
%\end{itemize}

\section{Discussão}

%Esta parte poderia ser integrada à seção de resultados, permitindo uma discussão mais direta. Por exemplo, suponha que você fez um ajuste de um modelo teórico aos dados experimentais. Este ajuste  pode ser apresentado diretamente com o gráfico da seção de "Resultados", porém sua discussão ficaria pendende.

\section{Conclusão}

% Serão valorizadas as discussões que agreguem uma interpretação física às observações.

%% Exemplos
\begin{figure}[H]
\centering
\includegraphics[width=0.3\textwidth]{frog.jpg}
\caption{\label{fig:frog}This frog was uploaded via the project menu.}
\end{figure}

Exemplo de referência à figura/tabela: \ref{fig:frog}

Exemplo de citação: \cite{greenwade93}.
% "project -> aaabib.bib" para editar o arquivo de referências. 

%% \section{Referências}
%Esta seção inclui as referências bibliográficas utilizadas no texto. Não economize aqui, coloque todas as referências utilizadas. No \LaTeX isto é muito fácil através do arquivo de referência .bib, o qual se usa com algum programa de gerência de bibliografias, como o Mendeley, Zotero, etc. Estes mesmos programas também são compatíveis com Microsoft Word, OpenOffice, etc.
\bibliographystyle{abntex2-num}
\bibliography{aaabib}

\end{document}

%% Exemplo de figura, copie e cole, substituindo os conteúdos.
\begin{figure}[H]
\centering
\includegraphics[width=0.3\textwidth]{frog.jpg}
\caption{\label{fig:frog}This frog was uploaded via the project menu.}
\end{figure}

%% Exemplo de tabela, copie e cole, substituindo os conteúdos.
\begin{table}[H]
\centering
\begin{tabular}{l|r}
Item & Quantity \\\hline
Widgets & 42 \\
Gadgets & 13
\end{tabular}
\caption{\label{tab:widgets}An example table.}
\end{table}